\documentclass[12pt,a4paper]{article}

% =========================================
% PAGE SETUP
% =========================================
\usepackage[
    a4paper,
    margin=25.4mm
]{geometry}

% =========================================
% BASIC PACKAGES
% =========================================
\usepackage{graphicx}       % Images
\usepackage{amsmath}        % Mathematics
\usepackage{amssymb}        % Mathematical symbols
\usepackage{booktabs}       % Professional tables
\usepackage{float}          % [H] figure/table positioning
\usepackage{caption}        % Caption customization
\usepackage{ragged2e}       % Text alignment
\usepackage{setspace}       % Line spacing
\usepackage{algorithm}      % Algorithms environment
\usepackage{algpseudocode}  % Pseudocode styling
\usepackage{hyperref}       % Clickable links and references

% =========================================
% APA REFERENCING
% =========================================
\usepackage[
    style=apa,
    backend=biber
]{biblatex}

\addbibresource{references.bib}

% =========================================
% HYPERLINK SETTINGS
% =========================================
\hypersetup{
    colorlinks=true,
    linkcolor=black,
    urlcolor=black,
    citecolor=black,
    pdftitle={Research Proposal - Assignment 2}
}

% Table and figure caption formatting
\DeclareCaptionFont{captionnine}{\fontsize{9pt}{11pt}\selectfont}

\captionsetup[table]{
    font=captionnine,
    justification=raggedright,
    singlelinecheck=false,
    labelsep=colon
}

\captionsetup[figure]{
    font=captionnine,
    justification=centering,
    singlelinecheck=false,
    labelsep=colon
}

% =========================================
% DOCUMENT
% =========================================
\begin{document}

\begin{center}

{\fontsize{18}{22}\selectfont\textbf{RESEARCH PROPOSAL}}

\vspace{0.5cm}

{\fontsize{16}{20}\selectfont\textbf{Dynamic Volume Optimisation And Adaptive Stopping Criteria}}

\vspace{1cm}

{\fontsize{12}{14}\selectfont\textbf{Group Members}}

\vspace{0.3cm}

{\fontsize{12}{14}\selectfont
1. [REDACTED] -- [REDACTED] - 33\% contribution \\
2. [REDACTED] -- [REDACTED] - 33\% contribution\\
3. [REDACTED] -- [REDACTED] - 33\% contribution\\
}

\vspace{0.8cm}

{\fontsize{12}{14}\selectfont
\textbf{Programme:} BACHELOR OF COMPUTER SCIENCE (HONOURS) \\
\textbf{Course:} CPT6123-RESEARCH METHOD IN COMPUTER SC\\
\textbf{Lecturer:} Ts. [REDACTED]\\
\textbf{Semester:} Trimester July/August 2025\\
\textbf{Submission Date:} 27th September 2026
}

\end{center}

% =========================================
% ABSTRACT
% =========================================
\vspace{0.5cm}

\noindent
\textbf{Abstract}

\vspace{0.2cm}

\begin{justify}
[Provide a concise overview of the proposed project (150--250 words). Summarize the technical problem, research objectives, proposed methodology/technical approach, and anticipated outcomes.]

\vspace{0.3cm}

\textbf{Keywords:} 4--6 technical keywords related to your project.
\end{justify}

% =========================================
% TABLE OF CONTENTS
% =========================================
\newpage
\tableofcontents
\newpage

% =========================================
% MAIN CONTENT
% =========================================
\begin{justify}

% =========================================
% 1. INTRODUCTION
% =========================================
\section{Introduction}

\subsection{Background and Context}
Introduce the broader Computer Science domain and context. Discuss the importance and relevance of the topic in recent years. Provide citations to support key claims.

\subsection{Motivation and Significance}
Explain why this research problem is critical to solve now, dvetailing its practical relevance, real-world utility, or industry significance.


% =========================================
% 2. LITERATURE REVIEW RECAP
% =========================================
\section{Literature Review Summary}
Briefly synthesize the findings from Assignment 1 to ground the proposal in existing research.

\begin{itemize}
    \item \textbf{Current State of Knowledge:} Overview of the primary techniques and frameworks currently used.
    \item \textbf{Identified Gaps from Assignment 1:} Explicitly outline the limitations and shortcomings in existing work that motivate this proposed solution.
\end{itemize}


% =========================================
% 3. PROBLEM STATEMENT
% =========================================
\section{Problem Statement}
Define the technical challenge using the \textbf{Context--Gap--Solution} format:
\begin{itemize}
    \item \textbf{Context:} The technical setting and current practices.
    \item \textbf{Gap:} The documented technical shortcoming or unresolved challenge.
    \item \textbf{Proposed Solution:} The high-level proposal designed to address this shortcoming.
\end{itemize}


% =========================================
% 4. RESEARCH QUESTIONS AND OBJECTIVES
% =========================================
\section{Research Questions and Research Objectives}

\subsection{Research Questions}

\begin{itemize}
    \item \textbf{RQ1:} How can an adaptive stopping criterion be developed to monitor real-time distribution shift and coverage saturation during synthetic sample generation for industrial soft sensors?
    \item \textbf{RQ2:} What is the performance impact of sensor accuracy when we automatically control the amount of generated data to prevent the data quality from getting worse when exceeding optimal sample thresholds (e.g., $\sim 420$ samples)?
    \item \textbf{RQ3:} How effectively does the proposed adaptive stopping framework balance computational speed and prediction accuracy compared to traditional, fixed-limit methods across different datasets?
    \item \textbf{RQ4:} To what extent does automated generation control help reduce unnecessary repeated data and signal noise in systems with limited computing power?
\end{itemize}

\subsection{Research Objectives}

\begin{itemize}
    \item \textbf{RO1:} To design an adaptive stopping criterion and a real-time monitor to track data pattern changes and coverage levels within the first 6 weeks of the project timeline.
    \item \textbf{RO2:} To evaluate and compare how controlling data volume helps prevent data distortion and keeps sensor accuracy high compared to fixed limits (beyond $\sim 420$ samples) over a 4-week testing phase.
    \item \textbf{RO3:} To check the reliability, speed, and prediction accuracy of the new framework against standard baseline methods using both standard test functions and industrial datasets within the final 4 weeks.
    \item \textbf{RO4:} To measure how much unnecessary data and noise are reduced when the automated generation control is used in simulated working environments.
\end{itemize}


% =========================================
% 5. RESEARCH METHODOLOGY
% =========================================
\section{Research Methodology}

\subsection{Research Design}
% Specify and justify the research design approach (e.g., \textit{Experimental, Empirical, Quasi-experimental, or Comparative}). Provide an architectural flow overview describing how data flows through the proposed pipeline.
Using an \textbf{experimental} approach, this study aims to determine whether the amount of synthetic data generated by a GAN affects the performance of downstream machine learning models, the similarity between real and synthetic data, and the computational cost of data generation. The main variable being investigated is therefore the volume of generated synthetic data and the point at which generation is stopped.\\

\noindent The experimental process is divided into three main phases:

\begin{enumerate}
    \item \textbf{Phase 1: Data Preparation} \\
    \noindent Datasets are preprocessed and converted into suitable feature representations.
    \item \textbf{Phase 2: Generation and Monitoring} \\ 
    \noindent GAN models are trained and used to generate synthetic samples in batches. A proposed ``Adaptive Comparing and Stopping'' framework then monitors the generated data after each batch and determines whether generation should continue or stop.
    \item \textbf{Phase 3: Evaluation and Comparison} \\
    The final augmented datasets are used to train downstream classifiers. Their classification performance, distribution quality, and computational cost are then compared with the baseline methods.
\end{enumerate}

\subsection{Technical Approach}
% Detail the proposed algorithms, network architectures, framework adaptations, or mathematical formulations. Distinguish clearly between existing frameworks repurposed and the novel components introduced.

\subsubsection{Proposed ``Adaptive Comparing and Stopping'' Framework}
This is the main framework for this proposed research methodology. In a typical approach, researchers decide in advance how many synthetic samples should be generated. This can result in either generating too few samples or generating more samples than necessary. \\

\noindent This framework takes a different approach. It generates synthetic data one batch at a time and checks the effect of each new batch before deciding whether to continue. \\

\noindent The basic process is: \\

Generate $\rightarrow$ Measure $\rightarrow$ Decide $\rightarrow$ Continue or Stop \\

\noindent Hence, this framework acts as a monitoring and stopping layer for the GAN. It compares the generated data with the original real data and looks for two situations to tell the GAN when to stop generating:

\begin{itemize}
    \item Saturation: additional samples are producing very little change.
    \item Distribution drift: additional samples are causing the synthetic distribution to move further away from the real distribution.
\end{itemize}

\noindent The overall system architecture is described below:

\begin{figure}[htbp]
    \centering
    \includegraphics[width=0.8\textwidth]{overallSystemArch.png}
    \caption{Overall Architecture of the Proposed ``Adaptive Comparing and Stopping'' Framework}
    \label{fig:overallSystemArch}
\end{figure}

\subsubsection{Saturation Detection}
The saturation detection works by monitoring how much the measured difference between the real and synthetic data changes as more batches are generated. When the change becomes very small, generating additional samples may provide little further benefit.
For example:
\begin{enumerate}
    \item Batch 1: Large change (continue)
    \item Batch 2: Large change (continue)
    \item Batch 3: Moderate change (continue)
    \item Batch 4: Small change (continue)
    \item Batch 5: Very small change (stop)
\end{enumerate}
When the difference remains below a defined saturation threshold, the framework considers that the generation process has reached a point of diminishing returns and stops generating more samples.

\subsubsection{Distribution Drift Detection}
Similarly, the proposed framework also includes a second stopping mechanism to detect possible distribution distortion on every batch. When the condition is detected:
\begin{enumerate}
    \item The most recent batch is removed.
    \item The generation process is stopped.
    \item The remaining synthetic data is returned as the final augmented dataset.
\end{enumerate}

\subsubsection{Pseudocode of “Adaptive Comparing and Stopping” Framework}
INPUT, real dataset, trained GAN, batch size, saturation threshold \\
SET synthetic\_dataset = empty \\
SET history = empty \\
SET stop = FALSE \\
EXTRACT features from real dataset \\
WHILE stop = FALSE \\
\indent GENERATE a batch of synthetic data \\
\indent ADD batch to synthetic\_dataset \\
\indent EXTRACT features from synthetic\_dataset \\
\indent COMPARE synthetic data with real data \\
\indent STORE the comparison result in history \\
\indent IF the data has reached saturation THEN \\
\indent \indent SET stop = TRUE \\
\indent ELSE IF the data shows distribution drift THEN \\
\indent \indent REMOVE the latest batch \\
\indent \indent SET stop = TRUE \\
\indent END IF \\
END WHILE \\
RETURN synthetic\_dataset \\

\subsubsection{Justification for the Proposed Approach}
Unlike conventional augmentation methods that rely on fixed rules or fixed sample limits, our proposed ``Adaptive Comparing and Stopping'' framework uses a feedback-based stopping mechanism that monitors generated data directly to decide when to stop. Previous work showed performance degraded after 420 synthetic samples, proving that more data isn't always better. This framework tests whether the optimal stopping point can be detected automatically from data behavior.

\subsection{Data and Tools}
\begin{itemize}
    \item \textbf{Datasets:} Specify public benchmarks or data sources (e.g., ImageNet, COCO, Kaggle, IEEE DataPort). Justify the data choice (size, noise, features) and state any ethical or privacy considerations.
    \item \textbf{Tools and Hardware:} List the development frameworks, libraries (e.g., PyTorch, TensorFlow, Wireshark), and experimental environments (GPUs, OS, compute resources).
\end{itemize}

\subsection{Evaluation Criteria}
Define quantifiable performance metrics and evaluation criteria:
\begin{itemize}
    \item \textbf{Quantitative Metrics:} (e.g., F1-score, accuracy, throughput, latency in ms, energy consumption in Joules).
    \item \textbf{Baselines and Ablation Study:} Describe benchmark comparison models and planned ablation studies to evaluate individual components.
\end{itemize}


% =========================================
% 6. EXPECTED OUTCOMES
% =========================================
\section{Expected Outcomes}
Detail the anticipated results, potential deliverables (e.g., source code, modular toolkit, trained model weights), and the potential scientific or practical contributions of this project.


% =========================================
% 7. RESEARCH ACTIVITIES AND MILESTONES
% =========================================
\section{Project Timeline}
Outline the phased roadmap, deliverables, and critical milestones across the execution period.

\begin{table}[H]
\caption{Proposed project execution timeline and milestones}
\label{tab:timeline}
\fontsize{10pt}{12pt}\selectfont
\begin{tabular}{|p{2.5cm}|p{6.5cm}|p{5cm}|}
\hline
\textbf{Phase / Week} & \textbf{Key Activities} & \textbf{Deliverables / Milestones} \\
\hline
Phase 1 (Weeks 1--2) & Literature refinement \& dataset acquisition & Cleaned dataset \& pre-processing scripts \\
\hline
Phase 2 (Weeks 3--5) & Model architecture design \& pipeline setup & Base system architecture established \\
\hline
Phase 3 (Weeks 6--8) & Model implementation \& parameter tuning & Algorithm implementation \\
\hline
Phase 4 (Weeks 9--11) & Comparative benchmarking \& ablation tests & Empirical results \& evaluation tables \\
\hline
Phase 5 (Weeks 12--14) & Documentation, analysis, \& final thesis drafting & Complete project report \\
\hline
\end{tabular}
\end{table}


% =========================================
% 8. CONCLUSION
% =========================================
\section{Conclusion}

\subsection{Summary of Contributions}
Summarize the problem, proposed solution, and intended methodology.

\subsection{Final Remarks}
Present final remarks highlighting the long-term impact and relevance of the proposed project.

\end{justify}

% =========================================
% REFERENCES
% =========================================
\newpage
\phantomsection
\addcontentsline{toc}{section}{References}
\printbibliography[title={References}]

\end{document}
